% ============================== УВОД ==============================
\chapter*{Увод}
\addcontentsline{toc}{chapter}{Увод}
\markboth{Увод}{}

\section*{Контекст}
Центровете за данни и публичните облаци обслужват уеб приложения, бази от данни, системи за обработка на големи данни и услугите за Интернет на нещата. Концентрацията на ресурси ги прави привлекателна цел. Атаките Distributed Denial of Service (DDoS) достигат терабитови мощности. Botnet Mirai през 2016~г. използва стотици хиляди заразени устройства за атаки срещу доставчици на DNS и хостинг услуги~\cite{antonakakis2017}. Едновременно с това сканирането на портове и автоматизираните опити за отгатване на пароли срещу публично достъпни услуги като SSH и FTP са постоянен фонов шум във всяка облачна инфраструктура. Сканирането е първата стъпка на почти всяка атака: с него нападателят открива кои услуги са достъпни~\cite{lyon2009}.

Защитата на центъра за данни не може да разчита само на агенти върху крайните точки. Много системи са виртуални машини или контейнери с кратък живот, а част от трафика минава между вътрешни услуги. Затова основна линия на защита е \emph{мрежовата система за откриване на прониквания} (Network Intrusion Detection System, NIDS), която наблюдава потоците и сигнализира при подозрително поведение. Сигнатурните NIDS, като Snort и Suricata, откриват само вече описани атаки. Подходите с машинно обучение (ML) се учат от примери и могат да разпознават закономерности, които не са описани ръчно~\cite{sommer2010,buczak2016}. Те обаче имат собствени рискове: зависят от качеството на данните за обучение, могат да постигат подвеждащо високи метрики при нереалистична оценка и могат да се преобучат.

\section*{Мотивация и изследователски въпроси}
Литературата за ML-базирани IDS често съобщава точност над 99\%. Такива резултати обаче рядко отчитат три фактора: (1) дали признаците на отделния запис изобщо носят отпечатъка на атаката, (2) реалната честота на атаките в трафика и (3) статистическата несигурност и стабилността на самите метрики. Настоящото изследване поставя тези фактори в центъра на експеримента и търси отговор на следните въпроси:

\begin{enumerate}[label=\textbf{В\arabic*.},leftmargin=2.4em]
  \item Колко точно три принципно различни алгоритъма (ансамбъл с bagging, метод с ядро и ансамбъл с boosting) откриват три атаки с различен характер в трафика на център за данни?
  \item Как редукцията на вектора на признаците (подбор на признаци или PCA) влияе на качеството и на скоростта на моделите?
  \item Статистически значими ли са разликите между алгоритмите, колко стабилни са метриките и има ли преобучение?
  \item Какво е реалното качество на моделите при естественото съотношение между атаки и нормален трафик и как може да се избере прагът на решение?
  \item Кои признаци на потока носят информацията за всяка атака и от какво зависи постижимата точност?
\end{enumerate}

\section*{Принос}
Работата има следния принос:
\begin{itemize}
  \item напълно възпроизводим експериментален конвейер с отворен код, който използва \emph{само реални} записи от CSE-CIC-IDS2018 и CIC-IDS2017 и предпазва от изтичане на информация между обучение и тест;
  \item обосновка на избора на набори от данни: предварително изследване върху CICIoT2023 показва количествено защо Port Scan там не може да се открие с точност над около 85\%, а същата атака в потоковото представяне на CIC-IDS2017 се открива с над 99\%;
  \item декларативна архитектура (регистри на атаки и алгоритми), при която нова атака, нов алгоритъм или атака от друг набор се добавят с един запис;
  \item девет модела, оценени по TP, TN, FP, FN, Accuracy, Precision, Recall, F1 и Specificity с тълкуване „на 100 случая“, 95\% доверителни интервали, тест на McNemar за всяка двойка алгоритми и явна диагностика на преобучението;
  \item пет контролни експеримента (E1--E5) и девет разширени анализа (A1--A9), сред които криви на обучение, сравнение на подбор на признаци и PCA, избор на праг при естествено съотношение и пренос между атаките;
  \item пет проверки за приложимост (A10--A14): без признаците-отпечатъци на инструмента, хронологично разделяне, пренос между подвидовете на атаката, нормален трафик от непознат ден и комбиниран детектор XGBoost или SVM.
\end{itemize}

\section*{Структура на доклада}
Структурата следва зададения план. \Cref{ch:method} описва целта и методологията: избора на атаки и набор от данни, предварителната обработка, математическите основи на алгоритмите и метриките за оценка. \Cref{ch:framework} представя експерименталната технологична рамка, а \cref{ch:plan} описва плана на експериментите с фаза на машинно обучение и оперативна фаза. \Cref{ch:results} съдържа имплементацията и резултатите, а \cref{ch:analysis} ги анализира. Следват заключение и информационни източници (\cref{ch:refs}). Приложенията съдържат речник на признаците, пълните резултати от настройката, таблиците на анализите A3 и A4, фрагменти от кода, допълнителни фигури и инструкции за възпроизвеждане.
